% Section 12 — The catastrophe: Aphec, Ichabod, Dagon, the plagues, the
% Beth-shemesh crux, Cariathiarim. Drafted from ot-corpus.md §§26–29,
% journey.md §§1.6–1.9, w2-textcrit.md §§2, 4, and w2-rights.md §§1.11–1.13.
\sectionguard
\section{The catastrophe: from Aphec to Cariathiarim}

Begin by giving the elders of Israel their due. Their army has just been
beaten at Aphec, and they ask the right question: why has the Lord struck us
today before the Philistines? (1 Kings 4:3; in the Douay-Rheims numbering
this study uses, 1 Kings 4 = 1 Samuel 4). Their answer is not vulgar
superstition; it is memory. Let the ark be fetched from Silo, they reason,
that it may save us out of the hand of our enemies. And the memory is
accurate. The ark went first into the Jordan, and the river stood; it
circled Jericho, and the wall came down. If the throne of the Lord goes
before the army, how can the army fall? The reader who has come this far
could sign that argument himself. It is exactly wrong. No prophet steps
forward to say so; the event itself will explain why. The precedents the
elders remember were never demonstrations of a lever that works when pulled.
They were acts of a free Lord. At Hormah, on another day, that same Lord
stayed in the camp while a presumptuous Israel marched without him and was
cut to pieces (Numbers 14:44). The elders have remembered the wonders and
forgotten the freedom. What they fetch from Silo, they think, is a weapon.
What arrives is a judge.

\subsection{Aphec: the glory departs}

So the people send to Silo, and there comes into the camp ``the ark of the
covenant of the Lord of hosts sitting upon the cherubims: and the two sons
of Heli, Ophni and Phinees, were with the ark of the covenant of
God''\endnote{Douay-Rheims (Challoner), 1 Kings 4:4, read at
\url{https://drbo.org/chapter/09004.htm}. Douay's 1--2 Kings are the books
called 1--2 Samuel in Hebrew-numbered Bibles; its 3--4 Kings are 1--2 Kings.
The section quotes the Douay text as printed throughout.} (1 Kings 4:4).
Pause on the English. The participle attaches oddly, so that on the surface
it is the ark that sits upon the cherubims. The phrase in fact renders the
great divine title --- the Lord enthroned upon the cherubim, the title that
binds the throne to the box --- and the study quotes it as printed and reads
it so. The verse's second clause is as ominous as its first is majestic.
Ophni and Phinees are the two priests the reader left at Silo: the sons of
Heli whose contempt for the offerings of the Lord has already been narrated
and sentenced (1 Kings 2). The ark comes up to Aphec carried, so to speak,
by its own indictment.

The camp of Israel shouts until the earth rings. The Philistines, learning
what has arrived, are afraid, remembering what these gods did to Egypt
(4:5--8). Israel remembers Jordan, Philistia the plagues: both armies hold
the elders' theology. Then the battle: ``And the ark of God was taken: and
the two sons of Heli, Ophni and Phinees, were
slain''\endnote{Douay-Rheims, 1 Kings 4:11, at
\url{https://drbo.org/chapter/09004.htm}.} (4:11). A man of Benjamin runs
all day to Silo with his clothes rent. Heli, ninety-eight years old, sits
waiting by the gate, and the messenger's report climbs from bad to worse:
Israel fled; a great slaughter; your two sons dead; and the ark of God is
taken. It is at the ark, not at his sons, that the old priest falls
backward and dies (4:17--18). The narrator has graded the losses, and a
father's grief comes second.

The last word of the chapter belongs to a dying woman. The wife of Phinees,
hearing that the ark is taken and her husband dead, is brought to labor
before her time, and with her last strength she names the child:

\begin{witnessbox}{1 Kings 4:21--22, Douay-Rheims}
And she called the child Ichabod, saying: The glory is gone from Israel,
because the ark of God was taken\ldots\ The glory is departed from Israel,
because the ark of God was taken.
\end{witnessbox}

The name is a theology compressed to three syllables.\endnote{Douay-Rheims,
1 Kings 4:21--22, at \url{https://drbo.org/chapter/09004.htm}; verse 21 is
elided in the witness block as shown (its remaining clauses, on her father
in law and her husband, were not re-verified verbatim for this study's
records).} It carries the Hebrew word \term{kabod}, glory --- the glory of
the visible Presence. And the text itself, twice over, makes the equation
the rest of this study will need: the ark taken \emph{is} the glory
departed. Israel's own Scripture never treats the loss at Aphec as the loss
of a treasured object. It treats it as an eclipse. Panel II will follow the
word \term{kabod} from this deathbed through the Psalter and the prophets.
For now it is enough that the equation is no later devotion: it is the
narrative's own first-generation verdict, put in the mouth of a woman with
nothing left to lose. The glory is departed. That is what the box meant.

\subsection{Dagon on his face}

``Now the ark of God was in the land of the Philistines seven
months''\endnote{Douay-Rheims, 1 Kings 6:1, at
\url{https://drbo.org/chapter/09006.htm}.} (1 Kings 6:1). Those seven
months are among the strangest campaigns in the Old Testament, because
Israel takes no part in them. The Philistines do what any ancient victor
did with a captured god: they carry the trophy into the house of their own
god, to declare whose deity won. The trophy declines the role. In the
morning at Azotus --- Douay's name for Ashdod --- ``behold Dagon lay upon
his face on the ground before the ark of the Lord'' (5:3). They set him up
again, and the second morning ``the head of Dagon, and both the palms of
his hands were cut off upon the
threshold''\endnote{Douay-Rheims, 1 Kings 5:3--4, at
\url{https://drbo.org/chapter/09005.htm}.} (5:4). The idol lies prostrate
like a courtier, then dismembered like a defeated king --- head and hands,
before the thing he was supposed to have conquered.

Then the plagues. The Douay prints at 5:6 a double affliction: the Lord
``afflicted Azotus\ldots\ with emerods. And in the villages and
fields\ldots\ there came forth a multitude of mice, and there was the
confusion of a great mortality in the city''\endnote{Douay-Rheims, 1 Kings
5:6, at \url{https://drbo.org/chapter/09005.htm}. ``Emerods'' is the old
English for tumors. The mice clause is a genuine version divergence: the
Vulgate, with the Septuagint, reads a plague of mice already at 5:6, where
the Masoretic Hebrew (so the Revised Standard Version and most moderns) has
only the tumors, mice first appearing at 6:4--5 as the golden images of the
guilt offering. The study prints its base text and reports the divergence;
it does not harmonize.} (5:6). The mice are a genuine version divergence
--- Vulgate and Septuagint against a Hebrew that keeps only the emerods ---
and the note discloses it. Either way the shape of the episode is the same.
The ark is passed like a hot coal from Azotus to Gath to Ekron, plague
breaking out at each remove (5:8--12), until the lords of the Philistines
are begged by their own people to send it home. No Israelite raid, no
ransom, no negotiation: the ark's captivity is conducted, from first to
last, as the Lord's solo campaign.

Even the sending home is made into evidence. The priests and diviners of
Philistia prescribe a guilt offering --- golden emerods and golden mice,
images of the plagues, five for the five lords (6:3--5). Then they design
what can only be called an experiment: a new cart, two milch kine that have
never borne the yoke, their calves shut up at home; ``if it go up by the way
of his own coasts towards Bethsames\ldots\ but if not, we shall know that it
is not his hand hath touched us, but it hath happened by
chance''\endnote{Douay-Rheims, 1 Kings 6:7--9, at
\url{https://drbo.org/chapter/09006.htm}; ``coasts'' in the older English
sense of borders. Bethsames is Douay's form of Beth-shemesh.} (6:7--9).
Untrained cattle, against every instinct, leaving their calves, taking the
straight road up to Israelite territory. The diviners have honestly left
the door open for coincidence, and the narrative lets pagan divination be
good science. The kine go straight up the highway to Bethsames, lowing as
they go, turning neither right nor left (6:12). The test convicts.

\subsection{Bethsames: the number as printed}

At Bethsames it is wheat harvest, and the reapers lift up their eyes and
rejoice; the cart is broken up for altar wood and the kine offered upon it
(6:13--14). The ark is home. Then, with no warning the narrative cares to
give, the day of joy becomes the section's hardest verse --- hardest in
what it reports, and hardest in how its text has come down. The number must
be quoted before it can be discussed:

\begin{witnessbox}{1 Kings 6:19, Douay-Rheims}
But he slew of the men of Bethsames, because they had seen the ark of the
Lord: and he slew of the people seventy men, and fifty thousand of the
common people. And the people lamented, because the Lord had smitten the
people with a great slaughter.
\end{witnessbox}

Fifty thousand and seventy dead in a harvest village.\endnote{Douay-Rheims,
1 Kings 6:19, at \url{https://drbo.org/chapter/09006.htm}, quoted in the
witness block as printed.} The number has troubled readers for as long as
the text has had readers. The Hebrew is itself the first problem: the
Masoretic text reads, word for word, ``of the people seventy men, fifty
thousand men'' --- two numbers set side by side with no connective, a
syntax anomalous in Hebrew counting and wildly out of scale for the place.
The Douay's ``seventy men, and fifty thousand of the common people''
follows the Vulgate's attempt to construe that anomaly. The Septuagint
keeps both figures and adds an offense of its own: the sons of Jechonias,
it says, did not rejoice when they saw the ark --- as though refusal of joy
were the crime.\endnote{Masoretic reading as printed in the Revised
Standard Version's translators' note at 1 Samuel 6:19: ``Cn: Heb \emph{of
the people seventy men, fifty thousand men}'' (RSV, read at
\url{https://www.biblegateway.com}); the RSV text prints ``he slew seventy
men of them'' --- a correction, so flagged by the translators themselves.
Septuagint (Rahlfs) at 1 Kingdoms 6:19, read via Blue Letter Bible
(blueletterbible.org): \term{kai ouk ēsmenisan hoi huioi Iechoniou en tois
andrasin Baithsamus, hoti eidan kibōton kyriou, kai epataxen en autois
hebdomēkonta andras kai pentēkonta chiliadas andrōn} --- the sons of
Jechonias did not rejoice among the men of Bethsames when they saw the ark
of the Lord; the clause is absent from the Hebrew, and the Greek keeps both
numbers. The standard apparatus further reports a few medieval Hebrew
manuscripts lacking the fifty thousand.} Josephus, retelling the scene in
the first century, gives seventy only, and his translator Whiston confessed
he could not account for the fifty thousand of the other
copies.\endnote{Josephus, \work{Antiquities} 6.1.4 (Niese §16), Greek text
read via the Scaife Viewer (scaife.perseus.org),
\url{urn:cts:greekLit:tlg0526.tlg001.perseus-grc2:6.16} --- the number
verified at the exact passage is \term{hebdomēkonta}, seventy. English from
William Whiston's translation, Project Gutenberg ebook 2848,
\url{https://www.gutenberg.org/ebooks/2848}: the wrath of God ``struck
seventy persons of the village of Bethshemesh dead, who, not being priests,
and so not worthy to touch the ark, had approached to it''; Whiston adds in
his note 4 that how other copies ``come to add such an incredible number as
fifty thousand in this one town, or small city, I know not.'' Note that
Josephus also supplies a cause the Hebrew does not --- approaching and
touching, not being priests --- an early harmonizing exegesis of the
offense.} Modern critical Bibles largely follow that instinct: the Revised
Standard Version prints ``seventy men'' and flags it, with its own honest
siglum, as a correction of the Hebrew. And one negative must be printed,
because the contrary claim circulates. The great Samuel scroll from Qumran,
4QSam-a, preserves verses on either side of 6:19 but not the verse itself,
and no other scroll has it. No Dead Sea manuscript reads ``seventy''; none
reads anything there at all.\endnote{Coverage per
\work{The Dead Sea Scrolls Bible} (Abegg, Flint, and Ulrich, 1999), 1 Samuel
apparatus, read in an archive.org ebook for this study's research records:
4QSam-a preserves 1 Kings (1 Samuel) 6:1--7, 12--13, 16--18, 20--21 --- not
verse 19. The primary edition is DJD XVII (2005), not consulted directly.}
The case for seventy rests on Josephus, on the shape of the Hebrew anomaly,
and on plausibility. That is a strong case, but a reconstruction, and this
study, which quotes its base text as printed, will not emend silently in
either direction.

The Douay tradition has its own reading of the verse, and it deserves to be
heard as reception rather than dismissed as apologetics. Challoner's note on
6:19 glosses ``seen'': ``And curiously looked into. It is likely this plague
reached to all the neighbouring country, as well as the city of
Bethsames.''\endnote{Challoner's note at 1 Kings 6:19, read on the same
drbo.org chapter page, quoted verbatim.} The second sentence stretches the
geography to carry the number. The first is the interesting one, because it
is not invented for the occasion. It wires Bethsames back to the transport
law of the wilderness and its ban on seeing the holy things unwrapped:
``Let not others by any curiosity
see the things that are in the sanctuary before they be wrapped up,
otherwise they shall die''\endnote{Douay-Rheims, Numbers 4:20, at
\url{https://drbo.org/chapter/04004.htm}; ``by any curiosity'' is the
Vulgate's interpretive gloss, matching Challoner's ``curiously looked into''
at 1 Kings 6:19 --- the version's own internal cross-reference, quoted here
as the Douay tradition's reading, not as a rendering of the Hebrew.}
(Numbers 4:20). On the Douay's own terms the men of Bethsames died of the
statute the porters of the ark had lived under since Sinai: not for malice,
but for treating the holy as a curiosity. Whether that is what the Hebrew's
bare ``because they had seen'' means, the philology cannot settle. What the
deaths mean --- what kind of thing holiness is, that joy and irreverence
can stand this close together --- is Panel II's question, and it will not
be rushed. The men of Bethsames themselves ask it first, in almost those
words: who shall be able to stand before this holy Lord? (6:20). Their
answer is the practical one. They send messengers up the hill country, to
ask someone else to take the ark.

\subsection{Cariathiarim: the long parking}

The men of Cariathiarim --- Kiriath-jearim, the wooded height above the
Jerusalem road --- come down and take it to ``the house of Abinadab in
Gabaa: and they sanctified Eleazar his son, to keep the ark of the
Lord''\endnote{Douay-Rheims, 1 Kings 7:1, at
\url{https://drbo.org/chapter/09007.htm}. ``Gabaa'' renders the Hebrew for
``the hill'' as a place name.} (1 Kings 7:1). Then the narrator closes the
whole catastrophe with a sentence of studied quiet: ``from the day the ark
of the Lord abode in Cariathiarim days were multiplied, (for it was now the
twentieth year,) and all the house of Israel rested following the
Lord''\endnote{Douay-Rheims, 1 Kings 7:2, at
\url{https://drbo.org/chapter/09007.htm}, with the parenthesis as
printed. The mood of the verse is itself a version divergence: the Vulgate's
\term{requievit} (``rested following the Lord'') stands against the Hebrew
as modern versions render it --- RSV ``lamented after the Lord'' --- rest in
one tradition, homesick mourning in the other. Quoted as printed;
divergence disclosed.} (7:2). Two honest difficulties sit in that sentence.
The first is arithmetic: the twenty years cannot span the interval the
story requires. The ark stays in Abinadab's house through Samuel's
judgeship and the whole reign of Saul, until David's men come for it ---
decades on any chronology. The old solutions are two, and both are still on
the table. Either the twenty years counts only from the ark's arrival to
the national repentance narrated immediately after it (7:3--14), or the
verse is the seam where an older ark story was closed and the Saul
narratives joined on, the arithmetic never reconciled. Report, not
resolution.

The second difficulty is mood, and it divides the witnesses. The Douay's
Israel ``rested following the Lord,'' while the Hebrew as the moderns read
it has Israel \emph{lamenting} after him --- a nation either at peace under
the parked Presence or homesick for it. The study quotes what its base text
prints and leaves both moods standing.

One frame from the scholarship earns its two sentences here. Miller and
Roberts read these chapters against the ancient Near Eastern literature of
captured gods, in which carrying off a god's image was never conceded to be
the god's defeat. The Aphec-to-Bethsames narrative, they argue, is Israel's
masterwork in that genre: the ark taken, and its God thereupon needing no
rescue, wrecking Dagon, scourging five cities, and coming home unescorted
and undefeated.\endnote{Patrick D. Miller, Jr., and
J. J. M. Roberts, \work{The Hand of the Lord} (Johns Hopkins University
Press, 1977) --- reported here from the study's research records at the level of the
book's thesis; the monograph was not itself opened for this study. It
stands within a longer source-critical debate over these chapters, running
since Leonhard Rost's 1926 monograph, whose existence the study notes
without adjudicating.} On that reading the seven months in Philistia are
not the interruption of the story but its thesis. Israel lost a battle; the
Lord lost nothing.

So the catastrophe ends not with a restoration but with a parking. The ark
does not go back to Silo; the reader will learn in time that there is no
Silo to go back to. Nor does it come to rest in any sanctuary. The old
tabernacle drifts separately through the story from now on, surfacing at
Nob and later at Gibeon, a tent without its throne (1 Kings 21; 3 Kings
3:4; 2 Paralipomenon 1:3). The ark simply waits, in a hill-town house, one
consecrated man keeping it, while the narrative turns its face away: to
Samuel's circuit, to the demand for a king, to Saul's rise and long ruin.
Days were multiplied. Twenty years and more the glory of Israel stands in
Abinadab's house at Gabaa, and a kingdom grows up below it. The next hand
laid on the ark will be a king's.
